\textbf{Phần 1. Trắc nghiệm lựa chọn}
\setcounterpageref{ex}{0}

\begin{ex} %[1D10V38]
Tính: $3 + 4 + 3$.
\choice
{$0$}
{\True $10$}
{$20$}
{$9$}
\loigiai{
Thực hiện: $3 + 4 = 7$; $7 + 3 = 10$.
\textbf{Đáp số:} $10$.\par
\textbf{Đáp án của câu hỏi là B}
}
\end{ex}

\begin{ex} %[1D10V38]
Tính: $10 - 4 - 3$.
\choice
{$2$}
{$4$}
{$1$}
{\True $3$}
\loigiai{
Thực hiện: $10 - 4 = 6$; $6 - 3 = 3$.
\textbf{Đáp số:} $3$.\par
\textbf{Đáp án của câu hỏi là D}
}
\end{ex}

\begin{ex} %[1D10V38]
Tính: $8 - 5 + 4$.
\choice
{\True $7$}
{$6$}
{$8$}
{$5$}
\loigiai{
Thực hiện: $8 - 5 = 3$; $3 + 4 = 7$.
\textbf{Đáp số:} $7$.\par
\textbf{Đáp án của câu hỏi là A}
}
\end{ex}

\begin{ex} %[1D10V38]
Điền số thích hợp: $6 + \dots = 10$.
\choice
{$3$}
{$5$}
{\True $4$}
{$2$}
\loigiai{
Số cần điền là: $10 - 6 = 4$.
\textbf{Đáp số:} $4$.\par
\textbf{Đáp án của câu hỏi là C}
}
\end{ex}

\begin{ex} %[1D10V38]
Điền số thích hợp: $9 - \dots = 2$.
\choice
{$6$}
{\True $7$}
{$8$}
{$5$}
\loigiai{
Số cần điền là: $9 - 2 = 7$.
\textbf{Đáp số:} $7$.\par
\textbf{Đáp án của câu hỏi là B}
}
\end{ex}

\begin{ex} %[1D10V38]
Tìm $x$ biết: $x + 2 = 5 + 4$.
\choice
{$6$}
{$8$}
{$5$}
{\True $7$}
\loigiai{
Vế phải: $5 + 4 = 9$. Ta có: $x + 2 = 9 \Rightarrow x = 9 - 2 = 7$.
\textbf{Đáp số:} $7$.\par
\textbf{Đáp án của câu hỏi là D}
}
\end{ex}

\begin{ex} %[1D10V38]
Tìm $y$ biết: $10 - y = 8 - 2$.
\choice
{\True $4$}
{$3$}
{$5$}
{$2$}
\loigiai{
Vế phải: $8 - 2 = 6$. Ta có: $10 - y = 6 \Rightarrow y = 10 - 6 = 4$.
\textbf{Đáp số:} $4$.\par
\textbf{Đáp án của câu hỏi là A}
}
\end{ex}

\begin{ex} %[1D10V38]
Số lớn nhất có một chữ số trừ đi số bé nhất có một chữ số bằng bao nhiêu?
\choice
{$8$}
{$10$}
{\True $9$}
{$7$}
\loigiai{
Ta có: $9 - 0 = 9$.
\textbf{Đáp số:} $9$.\par
\textbf{Đáp án của câu hỏi là C}
}
\end{ex}

\begin{ex} %[1D10V38]
Tổng của số chẵn lớn nhất có một chữ số và số lẻ bé nhất có một chữ số khác 0 là bao nhiêu?
\choice
{$8$}
{\True $9$}
{$10$}
{$7$}
\loigiai{
Số chẵn lớn nhất là 8, số lẻ bé nhất khác 0 là 1: $8 + 1 = 9$.
\textbf{Đáp số:} $9$.\par
\textbf{Đáp án của câu hỏi là B}
}
\end{ex}

\begin{ex} %[1D10V38]
Lan có 5 cái kẹo, mẹ cho thêm 3 cái. Sau đó Lan cho em 2 cái. Hỏi Lan còn lại bao nhiêu cái kẹo?
\choice
{$5$}
{$7$}
{$4$}
{\True $6$}
\loigiai{
Lan có: $5 + 3 - 2 = 6$ (cái kẹo).
\textbf{Đáp số:} $6$.\par
\textbf{Đáp án của câu hỏi là D}
}
\end{ex}

\begin{ex} %[1D10V38]
Trên cành có 8 con chim. 3 con bay đi, rồi lại có 4 con chim khác bay đến. Hỏi trên cành có bao nhiêu con chim?
\choice
{\True $9$}
{$8$}
{$10$}
{$7$}
\loigiai{
Số chim trên cành là: $8 - 3 + 4 = 9$ (con chim).
\textbf{Đáp số:} $9$.\par
\textbf{Đáp án của câu hỏi là A}
}
\end{ex}

\begin{ex} %[1D10V38]
Trong bể có 10 con cá. Nam vớt ra 4 con nuôi ở bể nhỏ, sau đó bể nhỏ chết mất 1 con. Hỏi bể nhỏ còn lại bao nhiêu con cá?
\choice
{$2$}
{$4$}
{\True $3$}
{$1$}
\loigiai{
Bể nhỏ còn lại: $4 - 1 = 3$ (con cá).
\textbf{Đáp số:} $3$.\par
\textbf{Đáp án của câu hỏi là C}
}
\end{ex}

\textbf{Phần 2. Trắc nghiệm đúng sai}
\setcounterpageref{ex}{0}

\begin{ex} %[1D10V38]
Xét tính đúng hoặc sai của các khẳng định toán học sau đây:
\choiceTF
{Kết quả của bài toán «Tìm số tự nhiên $a$ sao cho: $a + a = 8$. Giá trị của $a$ là bao nhiêu» là $5$.}
{\True Kết quả của bài toán «Tìm số tự nhiên $b$ sao cho: $b + b + b = 9$. Giá trị của $b$ là bao nhiêu» là $3$.}
{Có 10 que tính, bạn Bình xếp được 2 hình tam giác riêng biệt (mỗi hình 3 que). Bình còn thừa $5$ que tính.}
{\True Kết quả của bài toán «Số 10 gồm 7 và mấy» là $3$.}
\loigiai{
\begin{itemize}
    \item Ý a) Sai: Kết quả đúng là $4$, không phải $5$. Hai số giống nhau có tổng bằng 8 là 4 ($4 + 4 = 8$). Vậy $a = 4$.
\textbf{Đáp số:} $4$.
    \item Ý b) Đúng: Ba số giống nhau có tổng bằng 9 là 3 ($3 + 3 + 3 = 9$). Vậy $b = 3$.
\textbf{Đáp số:} $3$.
    \item Ý c) Sai: Kết quả đúng là $4$, không phải $5$. Bình dùng: $3 + 3 = 6$ (que tính).\nSố que tính còn thừa là: $10 - 6 = 4$ (que tính).
\textbf{Đáp số:} $4$.
    \item Ý d) Đúng: Số 10 gồm 7 và $3$ (vì $7 + 3 = 10$).
\textbf{Đáp số:} $3$.
\end{itemize}
\begin{itemize}
    \item a) \textbf{Sai}.
    \item b) \textbf{Đúng}.
    \item c) \textbf{Sai}.
    \item d) \textbf{Đúng}.
\end{itemize}
}
\end{ex}

\begin{ex} %[1D10V38]
Xét tính đúng hoặc sai của các bài toán và phép tính sau đây:
\choiceTF
{\True Kết quả của bài toán «Số 8 gồm 5 và mấy» là $3$.}
{\True Kết quả của bài toán «Hai số có tổng bằng 10 và hiệu bằng 4. Số lớn hơn là số mấy» là $7$.}
{Kết quả của bài toán «Hai số có tổng bằng 9 và hiệu bằng 1. Số bé hơn là số mấy» là $5$.}
{Có 10 que tính, bạn Bình xếp được 2 hình tam giác riêng biệt (mỗi hình 3 que). Bình còn thừa $5$ que tính.}
\loigiai{
\begin{itemize}
    \item Ý a) Đúng: Số 8 gồm 5 và $3$ (vì $5 + 3 = 8$).
\textbf{Đáp số:} $3$.
    \item Ý b) Đúng: Hai số có tổng bằng 10 và hiệu bằng 4 là 7 và 3. Số lớn hơn là số $7$.
\textbf{Đáp số:} $7$.
    \item Ý c) Sai: Kết quả đúng là $4$, không phải $5$. Hai số có tổng bằng 9 và hiệu bằng 1 là 5 và 4. Số bé hơn là số $4$.
\textbf{Đáp số:} $4$.
    \item Ý d) Sai: Kết quả đúng là $4$, không phải $5$. Bình dùng: $3 + 3 = 6$ (que tính).\nSố que tính còn thừa là: $10 - 6 = 4$ (que tính).
\textbf{Đáp số:} $4$.
\end{itemize}
\begin{itemize}
    \item a) \textbf{Đúng}.
    \item b) \textbf{Đúng}.
    \item c) \textbf{Sai}.
    \item d) \textbf{Sai}.
\end{itemize}
}
\end{ex}

\textbf{Phần 3. Trả lời ngắn}
\setcounterpageref{ex}{0}

\begin{ex} %[1D10V38]
Trong một hàng dọc có 8 bạn. Bạn đứng thứ tư có bao nhiêu bạn đứng phía sau mình?
\shortans{4}
\loigiai{
\textbf{Đáp số:} $4$\par
Số bạn đứng phía sau là: $8 - 4 = 4$ (bạn).
\textbf{Đáp số:} $4$.
}
\end{ex}

\begin{ex} %[1D10V38]
Tìm số đứng liền trước của số 9.
\shortans{8}
\loigiai{
\textbf{Đáp số:} $8$\par
Số liền trước số 9 là: $9 - 1 = 8$.
\textbf{Đáp số:} $8$.
}
\end{ex}

\begin{ex} %[1D10V38]
Tìm số đứng liền sau của số 7.
\shortans{8}
\loigiai{
\textbf{Đáp số:} $8$\par
Số liền sau số 7 là: $7 + 1 = 8$.
\textbf{Đáp số:} $8$.
}
\end{ex}

\begin{ex} %[1D10V38]
Có bao nhiêu số tự nhiên có một chữ số lớn hơn 6?
\shortans{3}
\loigiai{
\textbf{Đáp số:} $3$\par
Các số đó là: $7, 8, 9$. Có tất cả $3$ số.
\textbf{Đáp số:} $3$.
}
\end{ex}

\begin{ex} %[1D10V38]
Một hộp bút có 9 chiếc bút chì gồm hai màu đen và đỏ. Biết có 6 bút chì đen. Hỏi có bao nhiêu bút chì đỏ?
\shortans{3}
\loigiai{
\textbf{Đáp số:} $3$\par
Số bút chì đỏ là: $9 - 6 = 3$ (chiếc).
\textbf{Đáp số:} $3$.
}
\end{ex}

\begin{ex} %[1D10V38]
Tính: $10 - 1 - 2 - 3 - 4$.
\shortans{0}
\loigiai{
\textbf{Đáp số:} $0$\par
Trừ lần lượt: $10 - 1 = 9$; $9 - 2 = 7$; $7 - 3 = 4$; $4 - 4 = 0$.
\textbf{Đáp số:} $0$.
}
\end{ex}
